\documentclass[exercice,a4]{thedocument}%
\usepackage{tabularx}%
\usepackage[np]{numprint}%
\startdatakeys
    \source{25GENMATMEAG3}
    \BO{2026}
    \cycle{4}
    \competencesDuSocle{CA,RA}
    \objectifsApprentissage{C4A1aa,C4A2ad,C4B2cb,C4D2aa,C4D3ac}
    \motsCles{étendue,conversion,écriture scientifique,ratio,pourcentage,factorisation}
\enddatakeys
\begin{document}%
    \def\euro{\texteuro}%

\medskip

Cet exercice est un questionnaire à choix multiple(QCM). Pour chaque question, quatre
réponses (A, B, C ou D) sont proposées. Une seule réponse est exacte. Recopier sur la copie le numéro de la question et la lettre correspondant à la réponse exacte. Aucune justification n'est demandée.

\medskip

\begin{enumerate}
\item On considère la série suivante :

\[4~;~8~;~11~;~7~;~2~;~3~;~14\]

Quelle est l'étendue de cette série ?
\begin{center}
\begin{tabularx}{0.75\linewidth}{|*{4}{>{\centering \arraybackslash}X|}}\hline
A&B&C&D\\ \hline
10 &7& 12& 14\\ \hline
\end{tabularx}
\end{center}
\item Quel est le volume correspondant à 1 L ?

\begin{center}
\begin{tabularx}{0.75\linewidth}{|*{4}{>{\centering \arraybackslash}X|}}\hline
A&B&C&D\\ \hline 
1 m$^3$& 1 cm$^3$&1 dm$^3$&1 mm$^3$\rule[-2mm]{0mm}{6mm}\\ \hline
\end{tabularx}
\end{center}

\item Quel est le nombre dont l'écriture scientifique est $8,6 \times 10^{-4}$ ?

\begin{center}
\begin{tabularx}{0.75\linewidth}{|*{4}{>{\centering \arraybackslash}X|}}\hline
A&B&C&D\\ \hline 
\np{86000}& \np{0,00086}& \np{- 0,00086}& \np{0,000086}\\ \hline
\end{tabularx}
\end{center}

\item La longueur et la largeur du drapeau de la France sont dans le ratio $3~:~2$.

Quelle est la largeur du drapeau de la France dont la longueur est égale à 90~cm ? 

\begin{center}
\begin{tabularx}{0.75\linewidth}{|*{4}{>{\centering \arraybackslash}X|}}\hline
A&B&C&D\\ \hline
54 cm &135 cm& 45 cm&  60 cm\\ \hline
\end{tabularx}
\end{center}

\item Le prix d'un parfum est passé de $75$~\euro{} à $60$~\euro.

Quel pourcentage de réduction a été appliqué ?

\begin{center}
\begin{tabularx}{0.75\linewidth}{|*{4}{>{\centering \arraybackslash}X|}}\hline
A&B&C&D\\ \hline
80\,\%& 25\,\%& 15\,\%& 20\,\%\\ \hline
\end{tabularx}
\end{center}

\item Quelle est la forme factorisée de $4x^2 - 25$ ?

\begin{center}
\begin{tabularx}{0.75\linewidth}{|*{4}{>{\centering \arraybackslash}X|}}\hline
A			&B					&C					&D\\ \hline
$(2x - 5)^2$&$(2x - 5)(2x + 5)$& $(4x - 5)(4x + 5)$& $(4x - 5)^2$\rule[-2mm]{0mm}{6mm}\\ \hline
\end{tabularx}
\end{center}
\end{enumerate}

\bigskip
\showthesource
\end{document}
